%=============================================================================
\section{Specifying and Estimating Reliability}
\label{sec:reliability}
%=============================================================================

Equation~\eqref{eq:raig} requires $\hat\varepsilon_f$, which a deployed system
does not know. This section is about how little one can get away with knowing,
and it is the part of the argument on which the practical claim rests.

It is also the part a reader should press hardest, because the specification we
use throughout is three assumed constants. We therefore treat
$\hat{\boldsymbol\varepsilon}$ as something to be \emph{obtained} rather than
declared, and give five ways of obtaining it, in decreasing order of what they
demand of the practitioner: the truth (unattainable, an upper reference); a
three-tier curator annotation; a vector \emph{learned from the system's own
interaction logs}, which requires no annotation, no labels and no user study; a
vector computed from the catalogue's repeated items; and a single global
constant, which we show is provably inert. The third of these is the one that
answers the objection, and Section~\ref{sec:reliability-em} develops it.

\subsection{Five Specifications}
\label{sec:tiers}

\emph{Oracle.} $\hat\varepsilon_f = \varepsilon_f$. Not attainable; reported as
an upper reference.

\emph{Tiered.} Each attribute is annotated into one of $K$ qualitative
reliability tiers and assigned that tier's representative value. With $K=3$ the
annotation (Table~\ref{tab:tiers}) asks only whether an attribute is objective
(an unambiguous property
of the item, such as language of the lyrics), semi-objective (well defined but
imperfectly recalled, such as duration), or subjective (a perceptual judgement,
such as energy). This requires no measurement and can be produced by a domain
curator in minutes.

\emph{Learned from logs.} $\hat\varepsilon_f$ is estimated from the questions a
deployed system asked and the answers it received, by the procedure of
Section~\ref{sec:reliability-em}. No annotation, no labels, no user study, and
no knowledge of what any user was thinking of.

\emph{Empirical.} $\hat\varepsilon_f$ is measured from the catalogue itself by
the repeated-item estimator of Section~\ref{sec:reliability-empirical}. No
annotation, no users.

\emph{Uniform.} A single $\hat\varepsilon$ for all attributes. This is what a
system does when it models noise in the likelihood but not per attribute. It is
also, as the next proposition records, no correction at all.

\begin{proposition}[Uniform specification is inert for selection]
\label{prop:uniform}
If $\hat\varepsilon_f = \hat\varepsilon$ for every $f$, with
$\hat\varepsilon \in [0,\tfrac12)$, then
$\arg\max_q \RAIG(q) = \arg\max_q \Hb(p_q)$.
\end{proposition}

\begin{proof}
$\Hb(\hat\varepsilon)$ is then a constant additive offset, so it suffices that
$p \mapsto \Hb(p \star \hat\varepsilon)$ and $p \mapsto \Hb(p)$ induce the same
order. Both are symmetric about $p=\tfrac12$, so take $p_1 < p_2 \le \tfrac12$.
Since $p \star \hat\varepsilon = \hat\varepsilon + p(1-2\hat\varepsilon)$ is
strictly increasing in $p$ for $\hat\varepsilon < \tfrac12$, and maps
$[0,\tfrac12]$ into $[\hat\varepsilon,\tfrac12]$, we get
$p_1 \star \hat\varepsilon < p_2 \star \hat\varepsilon \le \tfrac12$, and $\Hb$
is strictly increasing on $[0,\tfrac12]$. Hence
$\Hb(p_1 \star \hat\varepsilon) < \Hb(p_2 \star \hat\varepsilon)$ exactly when
$\Hb(p_1) < \Hb(p_2)$.
\end{proof}

Proposition~\ref{prop:uniform} is worth stating plainly because it delimits the
contribution: heterogeneity is not a refinement of the method, it is the
entirety of the method. A system that models answer noise with one global rate
has corrected its belief update and left its selector exactly where it was.
This also predicts an experimental invariant we use as an internal
consistency check: RAIG with a uniform $\hat\varepsilon$ must reproduce
classical IG's accuracy trial for trial, which it does
(Section~\ref{sec:results-main}).

\begin{table}[t]
\caption{Attribute reliability annotation for the music catalogue. Crossover
values are assumed, informed by reported inter-annotator agreement on
perceptual music descriptors~\cite{hu2007mood,kim2010emotion}, and are not
measured from users. Section~\ref{sec:reliability-empirical} measures a related
quantity from the catalogue, and Section~\ref{sec:results-ensemble} reports the
primary comparison across a distribution over these values rather than at the
point values alone.}
\label{tab:tiers}
\centering
\footnotesize
\begin{tabular}{@{}lp{0.58\columnwidth}c@{}}
\toprule
\textbf{Tier} & \textbf{Attributes} & \textbf{$\varepsilon$} \\
\midrule
Objective & genre, language, content rating, artist type, release type, track version & 0.03 \\
Semi-objective & popularity, duration, liveness, instrumentalness, acousticness, loudness, speechiness & 0.15 \\
Subjective & mood, tempo, danceability, energy, valence, key, mode, time signature & 0.30 \\
\bottomrule
\end{tabular}
\end{table}

\subsection{Measuring Attribute Instability from Repeated Items}
\label{sec:reliability-empirical}

The most serious objection to the specification above is that its three numbers
are assumed. We cannot remove that objection without a user study, and
Appendix~\ref{app:protocol} specifies one. We can, however, measure something
adjacent, for free, in any catalogue with the right structure.

Many catalogues contain the same underlying entity more than once under
different identifiers: a song released on a single and again on an album, a
product listed by two vendors, a component under two part numbers. Those
duplicates are independent renderings of one item, so the rate at which their
\emph{discretised} attribute values disagree measures how fragile that
attribute's boundary is. An attribute whose two renderings of the same song
disagree nine per cent of the time has a boundary sitting inside the data,
which is precisely the property that makes it hard for a person to answer.
Formally, for attribute $f$ let $\mathcal{P}$ be the set of within-entity pairs;
the estimator is
\begin{equation}
d_f \;=\; \frac{1}{|\mathcal{P}|}\sum_{(i,j)\in\mathcal{P}}
          \mathbb{1}\!\left[a_f(c_i) \neq a_f(c_j)\right],
\label{eq:discordance}
\end{equation}
reported with a Wilson score interval~\cite{wilson1927}. This is the
repeated-observation logic of Dawid and Skene~\cite{dawid1979} applied to the
catalogue rather than to annotators.

In the music catalogue, $3{,}178$ (title, artist) groups contain two or more
distinct track identifiers, yielding $5{,}073$ within-group pairs.
Table~\ref{tab:discordance} reports the result.

\begin{table}[t]
\caption{Measured attribute instability $d_f$ from $5{,}073$ repeated-item
pairs, with Wilson 95\% intervals. This is a measurement, not an assumption.
It is not a user error rate: it is the rate at which the catalogue itself
disagrees with the catalogue about the same song.}
\label{tab:discordance}
\centering
\footnotesize
\begin{tabular}{@{}llcc@{}}
\toprule
\textbf{Attribute} & \textbf{Tier} & \textbf{$d_f$} & \textbf{95\% CI} \\
\midrule
language                & objective & 0.0000 & [0.0000, 0.0008] \\
artist type             & objective & 0.0000 & [0.0000, 0.0008] \\
release type            & objective & 0.0000 & [0.0000, 0.0008] \\
track version           & objective & 0.0000 & [0.0000, 0.0008] \\
content rating          & objective & 0.0158 & [0.0127, 0.0196] \\
genre                   & objective & 0.1234 & [0.1146, 0.1327] \\
\midrule
speechiness             & semi & 0.0041 & [0.0027, 0.0063] \\
instrumentalness        & semi & 0.0278 & [0.0236, 0.0327] \\
duration                & semi & 0.0375 & [0.0326, 0.0430] \\
acousticness            & semi & 0.0781 & [0.0710, 0.0858] \\
liveness                & semi & 0.0911 & [0.0835, 0.0993] \\
loudness                & semi & 0.1104 & [0.1021, 0.1193] \\
popularity              & semi & 0.4360 & [0.4224, 0.4497] \\
\midrule
time signature          & subjective & 0.0235 & [0.0196, 0.0280] \\
mode                    & subjective & 0.0467 & [0.0412, 0.0529] \\
key                     & subjective & 0.0603 & [0.0541, 0.0672] \\
tempo                   & subjective & 0.0645 & [0.0580, 0.0716] \\
danceability            & subjective & 0.0712 & [0.0644, 0.0786] \\
mood                    & subjective & 0.0785 & [0.0714, 0.0862] \\
valence                 & subjective & 0.0887 & [0.0812, 0.0968] \\
energy                  & subjective & 0.0928 & [0.0852, 0.1011] \\
\midrule
\multicolumn{2}{@{}l}{\emph{objective mean}}      & 0.0232 & \\
\multicolumn{2}{@{}l}{\emph{non-objective mean}}  & 0.0874 & \\
\bottomrule
\end{tabular}
\end{table}

The measurement partially corroborates the annotation and partially challenges
it, and both halves are informative. Objective attributes are $3.8\times$ more
stable than the rest ($0.0232$ against $0.0874$), which is the contrast the
two-tier ablation of Section~\ref{sec:results-ablation} depends on, and four of
the six objective attributes never disagree at all across five thousand pairs.
The rank correlation between the annotated $\varepsilon$ and the measured
$d_f$ is $\rho = 0.42$ ($p = 0.061$), positive as the annotation predicts but
not decisive.

The two attributes responsible for most of the disagreement are diagnosable
rather than mysterious, and neither is a perceptual judgement.
\emph{Popularity} is release-dependent by construction: a single and an album
cut of the same song genuinely have different popularity scores, so its
instability of $0.436$ measures a real property of the catalogue that has
nothing to do with whether a listener can answer the question. \emph{Genre} is
assigned per playlist rather than per track in this catalogue, so the same
recording appearing under two genre labels is a labelling convention, not a
boundary effect. Excluding those two, the correlation between annotation and
measurement rises to $\rho = 0.64$ ($p = 0.003$).

The honest reading is therefore: the estimator validates the objective versus
non-objective distinction strongly, does not by itself separate semi-objective
from subjective, and is confounded wherever an attribute varies across
renderings of an entity for reasons unrelated to perceptual vagueness. It is a
lower bound and a cross-check, not a replacement for measuring users.

\emph{As a specification.} We nonetheless evaluate it as one. RAIG-empirical
sets $\hat\varepsilon_f = \mathrm{clip}(\kappa d_f, 0.01, 0.45)$, where the
single global constant $\kappa$ is fixed so that the mean of
$\hat{\boldsymbol\varepsilon}$ matches the mean of the tier annotation
($\kappa = 2.50$). The shape then comes entirely from data and exactly one
number is assumed, against three for the tiered specification.
Section~\ref{sec:results-main} reports how it performs, and the answer is
instructive.

\subsection{Learning the Crossover from Interaction Logs}
\label{sec:reliability-em}

The tier values are assumed. The measurement of the previous subsection is free
but, as we reported, does not work well enough to deploy. This subsection gives
the specification that resolves the tension, and it requires nothing that a
running system does not already have.

\emph{Setting.} A system has been deployed with the incumbent selector---
classical information gain at a single global crossover, which is what a system
that has never heard of this paper would ship. It has logged, for each session
$d$, the questions it asked, $q_{d,1:T}$, and the answers it received,
$y_{d,1:T}$. It has \emph{not} logged what the user was thinking of, because it
does not know: the target item is latent. The question is whether
$\boldsymbol\varepsilon$ can be recovered from these logs alone.

It can, by the same argument Dawid and Skene use to recover annotator error
rates without ground truth~\cite{dawid1979}, with the item playing the role of
the latent label. Let $\bb_d$ be the posterior over items implied by session
$d$'s log under a current estimate $\hat{\boldsymbol\varepsilon}$. The
probability that the answer given at turn $t$ disagrees with the (unknown) true
item is then computable in closed form from the split mass
$p_{q} = \sum_c \bb_d(c)\,x_{q}(c)$:
\begin{equation}
\Pr\!\left[\,y_{d,t} \neq x_{q_{d,t}}(c^\ast)\,\right] =
\begin{cases} 1 - p_{q_{d,t}}, & y_{d,t} = 1,\\ p_{q_{d,t}}, & y_{d,t} = 0.\end{cases}
\label{eq:em-disagree}
\end{equation}
Averaging \eqref{eq:em-disagree} over every logged turn that probed attribute
$f$ gives the M-step,
\begin{equation}
\hat\varepsilon_f \;\leftarrow\;
\frac{1}{|\mathcal{T}_f|}\sum_{(d,t)\in\mathcal{T}_f}
  \Pr\!\left[\,y_{d,t} \neq x_{q_{d,t}}(c^\ast)\,\right],
\label{eq:em-mstep}
\end{equation}
and recomputing the posteriors under the updated estimate is the E-step.
Alternating the two is EM. The fixed point is a self-consistent account of the
logs: the crossover vector under which the answers actually received are as
likely as possible, given that nobody ever observed what anyone meant.

Two properties of this estimator matter more than its derivation.

\emph{It needs the ordering, not the level.} Proposition~\ref{prop:uniform}
says a constant added to every $\hat\varepsilon_f$ changes no selection.
An estimator that is systematically attenuated but correctly ordered is
therefore as good for selection as the truth, and this is not a hypothetical
concession: our estimates \emph{are} attenuated, systematically, on exactly the
tier where it would matter most (Section~\ref{sec:results-em}). That they work
anyway is a consequence of Proposition~\ref{prop:uniform}, not an accident.

\emph{Coverage is a design variable.} A greedy incumbent concentrates its
questions on a few attributes, so its logs may contain almost no observations
of the attributes it dislikes---which, by the bias this paper describes, are
precisely the reliable ones. A system can buy coverage by asking a small
fraction of its questions uniformly at random. We treat that fraction as a
parameter and sweep it, because whether it is necessary is an empirical
question rather than an obvious one, and the answer surprised us.

\subsection{A Distribution over the Tier Values}
\label{sec:priors}

Reporting a result at three assumed point values invites the question of
whether it survives different ones. We therefore place a prior on each tier and
report the primary comparison across draws from it
(Section~\ref{sec:results-ensemble}).

The prior is a scaled Beta on $[0,\tfrac12)$, not a Beta on $[0,1]$: the model
constrains $\varepsilon < \tfrac12$, since at $\tfrac12$ the answer is
independent of the item and beyond it the user is systematically
anti-correlated, which is a different phenomenon and not something a crossover
probability should represent. We draw
$\varepsilon = \tfrac12 X$ with $X \sim \mathrm{Beta}(2mc,\,(1-2m)c)$, giving
mean $m$ equal to the tier's point value and support where the model allows.
Concentrations are chosen low, since a tight prior would beg the question. The
resulting 95\% intervals are $[0.000,0.121]$ for objective,
$[0.037,0.300]$ for semi-objective and $[0.142,0.437]$ for subjective; note
that the semi-objective and subjective intervals overlap substantially, so the
ensemble includes many worlds in which the tier ordering the system assumes is
locally violated.

Crucially, the \emph{deployed} system is held fixed across the ensemble: it
always assumes the point values $(0.03, 0.15, 0.30)$, because that is what a
practitioner following this paper would ship. Only the truth varies, and it
varies per attribute, producing within-tier heterogeneity that the tiered
specification cannot represent. The question the ensemble answers is therefore
not whether the method works when it knows the answer, but whether a system
built on three guessed numbers still beats current practice when the world does
not match those numbers.

\subsection{Why Binning Makes Perceptual Attributes Fragile}
\label{sec:psychophysical}

Proposition~\ref{prop:discretisation} says binning fixes apparent
informativeness. A companion question is what binning does to
\emph{answerability}, and it admits a direct calculation. Suppose a listener
perceives a $z$-scored feature as $v + \sigma\zeta$ with
$\zeta \sim \mathcal{N}(0,1)$, and reports the quantile bin of that perceived
value. Because equal-population quantile binning necessarily places every
boundary in the densest region of the distribution, the mass lying within one
perceptual standard deviation of a boundary---and hence the probability of
reporting the wrong bin---is close to maximal by construction.

Evaluating this by Monte Carlo over the actual feature columns gives a bin-flip
probability, averaged over the ten continuous attributes whose raw column is
present in the shipped catalogue, of $0.077$ at $k=2$ and $\sigma=0.10$, rising
to $0.132$ at $\sigma=0.25$ and $0.200$ at $\sigma=0.50$. At $k=4$ the same
perceptual noise produces $0.165$, $0.288$ and $0.423$. The two mechanisms
therefore pull in the same direction as the discretiser is refined: finer
binning raises apparent informativeness towards $\Hb(1/k)$ while simultaneously
roughly doubling the error rate, which is exactly the trade a criterion blind
to $\varepsilon$ cannot see. It is also why the measured bias strengthens
rather than weakens under four-bin discretisation in
Section~\ref{sec:results-bias}.

This calculation is worth reading a second time as a partial defence of the
assumed tier values, because it is derived from the catalogue's own feature
distributions rather than borrowed. A listener with perceptual noise of a
quarter of a standard deviation, answering four-bin questions, misreports the
bin $0.288$ of the time on average---which brackets the subjective tier's
assumed $0.30$ from a direction entirely independent of the annotation. The
same listener answering two-bin questions misreports $0.132$ of the time,
against the semi-objective tier's assumed $0.15$. We do not present this as a
measurement of $\varepsilon$: it assumes a Gaussian perceptual noise level that
is itself unmeasured, and it says nothing about the objective tier, where the
error is one of recall rather than of perception. What it does establish is
that the assumed values are not arbitrary numbers chosen to make the method
work---they are what a standard psychophysical model produces on these
columns at a plausible noise level.
